\documentclass[./examen.tex]{subfiles}

\begin{document}
\begin{center}
\makebox[\textwidth]{Nombre y Apellido:\enspace\hrulefill}
\end{center}
\begin{questions}
    \question Un condensador que se conecta a una tensión de $V=200V$, adquiere cada armadura una carga de $Q=6x10^{-9}C$. Calcular su capacidad.
    \question La carga que adquiere un condensador al conectarlo a $V=900V$ es de $Q=0,007C$. Calcular su capacidad, y la capacidad equivalente si se conecta 3 iguales en serie.
       \question Calcular la permitividades absolutas para la $\epsilon_{nafta}=2.35 \quad, \epsilon_{aceite}=2.8\quad y \quad \epsilon_{agua}=80,5$.
   
    \question Un capacitor de 2 placas paralelas de $10cmx10cm$ si el dieléctrico es el vacío y la distancia entre placas de 0.1mm, calcular su capacidad.
    
    \question Una bobina de $N=400$ espiras es recorrida por una corriente $I=20A$ que da lugar a un flujo magnético de $\Phi=0,001Wb$. Calcular el coeficiente de autoinducción $L$ de la bobina.
    
\end{questions}


 \end{document}